\documentclass[11pt]{article}
\usepackage[utf8]{inputenc}
\usepackage[T1]{fontenc}
\usepackage{amsmath,amssymb}
\usepackage[margin=1.1in]{geometry}
\usepackage{booktabs}
\usepackage{authblk}
\usepackage[colorlinks=true,linkcolor=blue,citecolor=blue,urlcolor=blue]{hyperref}

\title{Ionization Energy Without Differencing:\\
\large One Constant for the Alkali Column}

\author[1]{Claes Johnson}
\affil[1]{KTH Royal Institute of Technology, SE-100 44 Stockholm, Sweden \\
\small\texttt{claesjohnson@gmail.com}}

\date{\today}

\begin{document}
\maketitle

\begin{abstract}
A first ionization energy is a small difference of large quantities: for potassium it is
$0.027\%$ of the total electronic energy. Any method whose total energies carry a per-cent
residual therefore cannot deliver it by differencing two solves --- the signal falls below
the error bar on each. This is why the standard formalism computes ionization by separate
machinery rather than by subtracting ground-state energies, and it is the reason a
parameter-free continuum model of the atom, accurate to one or two per cent on totals, is
silent on ionization if asked the question that way. We report a construction that does not
ask it that way. Screening the core to its net charge and giving the valence its own domain
outside a radius $R$ makes the core energy cancel \emph{analytically} rather than
numerically, leaving a one-electron problem whose energy \emph{is} the ionization energy.
At $R=0$ the construction is exact and unfittable: hydrogen returns $0.500$ and, on a kernel
of $+2$ with two valence electrons, helium returns $0.903$ against a measured $0.9036$. With
no outer cutoff the radius is the only length in the problem, and it is not free: against
Shannon's six-coordinate crystal radii a single constant, $R = 0.439\,r_{\mathrm{ion}}$,
reproduces the measured first ionization energies of lithium, sodium, potassium, rubidium
and caesium to within $4\%$, four of the five to $2\%$, and the result is robust to the
choice of radius convention. One parameter standing against five observables is a claim that
can fail. We state plainly what it is not: the radius is supplied from measurement rather
than computed, so this is a correlation between two measured quantities mediated by the
model, and gold does not obey it.
\end{abstract}

\section{The arithmetic that forbids differencing}\label{sec:arith}

Write $\mathrm{IE} = E(\mathrm{ion}) - E(\mathrm{atom})$. Both terms are large and nearly
equal. Table~\ref{tab:frac} gives the fraction of the total that the difference represents,
together with what a $1.6\%$ residual on the total --- the mean error of the continuum model
of~\cite{Johnson-atoms} across the first three periods --- amounts to in the same units.

\begin{table}[t]
\centering
\caption{A first ionization energy as a fraction of the total electronic energy it would have
to be differenced out of, and the signal-to-error-bar ratio at a $1.6\%$ residual on totals.
Measured values from~\cite{NIST}.}
\label{tab:frac}
\begin{tabular}{lrrrr}
\toprule
 & $|E_{\rm tot}|$ (Ha) & IE (Ha) & IE$/|E_{\rm tot}|$ & signal/error bar \\
\midrule
He & 2.903  & 0.9036 & $31.1\%$  & 19.4 \\
Li & 7.478  & 0.1981 & $2.65\%$  & 1.66 \\
Be & 14.667 & 0.3425 & $2.34\%$  & 1.46 \\
B  & 24.654 & 0.3049 & $1.24\%$  & 0.77 \\
C  & 37.845 & 0.4138 & $1.09\%$  & 0.68 \\
N  & 54.589 & 0.5341 & $0.98\%$  & 0.61 \\
O  & 75.067 & 0.5001 & $0.67\%$  & 0.42 \\
F  & 99.734 & 0.6403 & $0.64\%$  & 0.40 \\
Ne & 128.94 & 0.7925 & $0.62\%$  & 0.38 \\
Na & 162.25 & 0.1886 & $0.116\%$ & 0.07 \\
K  & 599.16 & 0.1595 & $0.027\%$ & 0.02 \\
\bottomrule
\end{tabular}
\end{table}

\textbf{From boron onward the quantity sought is smaller than the uncertainty on the
quantities it is extracted from.} At sodium it is one fourteenth of that uncertainty; at
potassium one fiftieth. The conclusion does not depend on which model supplies the totals.
It is arithmetic, and it applies to any method with a per-cent residual.

Two consequences follow. The first is that the standard formalism's practice is not a
historical accident: Koopmans' theorem, $\Delta$SCF, equation-of-motion coupled cluster and
$GW$ exist because ionization is computed by machinery built for it, not by subtracting two
ground-state energies~\cite{Koopmans,SzaboOstlund,Hedin}. The second is that a model
reporting totals to one or two per cent, and declining to report ionization energies, is not
thereby incomplete in a way the totals conceal --- it is declining to divide a number by
something smaller than its own error bar.

$\Delta$SCF nevertheless succeeds where the core is nearly unchanged between the two solves,
because then most of the error cancels between them. The construction below makes that
cancellation exact instead of accidental.

\section{The construction}\label{sec:constr}

Let the atom have $V$ valence electrons outside a closed core of $Z-V$ electrons. Replace
the core by its net charge, a kernel of charge $+V$, and let the valence densities occupy the
region outside a radius $R$. Writing $E_{\mathrm{val}}(n)$ for the energy of $n$ such
densities on that kernel,
\begin{equation}
\mathrm{IE} \;=\; E_{\mathrm{val}}(V-1) \;-\; E_{\mathrm{val}}(V),
\label{eq:ie}
\end{equation}
and the core energy is absent from both terms rather than subtracted between them. For
$V=1$ this reduces to $\mathrm{IE} = -E_{\mathrm{val}}(1)$: the ionization energy
\emph{is} the valence electron's energy, and no difference is formed at all.

The electron obeys
\begin{equation}
-\tfrac12\nabla^2 u \;-\; \frac{V}{r}\,u \;=\; \varepsilon\,u, \qquad r>R,
\label{eq:pde}
\end{equation}
with $u$ vanishing at $r=R$ and free at infinity, normalised to one electron. For $V>1$ the
valence densities are non-overlapping, so~\eqref{eq:ie} is a difference within one
composition, taken on one kernel, with the core absent from both sides.

\paragraph{Two checks in which nothing can be adjusted.} At $R=0$ the construction has no
freedom. For $V=1$ it is hydrogen and returns $-0.5$ exactly. For $V=2$ on a kernel of $+2$
it is helium, and~\eqref{eq:ie} returns $E(\mathrm{He}^+)-E(\mathrm{He}) = -2.000 -
(-2.9036) = 0.9036$; the computed value is $0.903$~\cite{Johnson-blogAE}. These are not fits.
They are the two points at which the radius is absent, and the construction meets them.

\paragraph{The outer boundary must be removed.} Earlier numerical realisations confined the
valence density to an annulus $R<r<5$~a.u.~\cite{Johnson-code1}, while the two-valence
realisation~\cite{Johnson-code2} limits it only by the box; the two are therefore not on the same
footing. That second length does real work: solved on
$R<r<\infty$ instead, the energy at $R=1.6$ changes from $0.196$ to $0.1447$~Ha. Since the
inner radius is the quantity of interest, the outer cutoff must go, and
with it the only other length in the problem. Equation~\eqref{eq:pde} on $R<r<\infty$ is
well posed, and we solve it as a radial eigenvalue problem converged to $10^{-8}$.

\section{One constant for the column}\label{sec:oneconst}

Table~\ref{tab:alkali} inverts~\eqref{eq:pde} for each alkali: the radius at which the
computed energy equals the measured ionization energy. The question is then whether those
radii are anything, or merely five numbers fitted to five numbers.

\begin{table}[t]
\centering
\caption{The radius required to reproduce each measured first ionization energy, against the
ion's measured six-coordinate radius~\cite{Shannon1976}. Crystal scale,
$r(\mathrm{O}^{2-})=126$~pm, which Shannon states corresponds more closely to the physical
size of ions in a solid.}
\label{tab:alkali}
\begin{tabular}{lrrrr}
\toprule
 & IE (Ha) & $R_{\rm fit}$ (a.u.) & $r_{\rm ion}$ (a.u.) & $R_{\rm fit}/r_{\rm ion}$ \\
\midrule
Li & 0.19814 & 0.8195 & 1.7008 & 0.482 \\
Na & 0.18886 & 0.9084 & 2.1921 & 0.414 \\
K  & 0.15952 & 1.2779 & 2.8724 & 0.445 \\
Rb & 0.15351 & 1.3756 & 3.1369 & 0.439 \\
Cs & 0.14310 & 1.5689 & 3.4204 & 0.459 \\
\midrule
\multicolumn{4}{l}{mean} & $0.448$ \\
\multicolumn{4}{l}{standard deviation} & $5.0\%$ \\
\bottomrule
\end{tabular}
\end{table}

They are not five free numbers. \textbf{The required radius is a fixed fraction of the ion's
measured radius.} Fixing that fraction once at $k=0.439$ and computing the column with no
per-atom adjustment gives
\begin{equation}
R = 0.439\,r_{\rm ion} \;\Longrightarrow\;
\frac{\mathrm{IE_{comp}}}{\mathrm{IE_{meas}}} =
1.04,\;\, 0.97,\;\, 1.01,\;\, 1.00,\;\, 1.02
\quad (\mathrm{Li},\mathrm{Na},\mathrm{K},\mathrm{Rb},\mathrm{Cs}).
\label{eq:oneconst}
\end{equation}
All five within $4\%$; four within $2\%$. One constant against five observables.

\paragraph{Robust to the convention.} Shannon tabulates two six-coordinate scales differing
by exactly $14$~pm: crystal radii, $r(\mathrm{O}^{2-})=126$~pm, and effective radii,
$r(\mathrm{O}^{2-})=140$~pm, the latter chosen for consistency with Pauling. On the effective
scale the constant becomes $k=0.484$ and the agreement is comparable --- within $8\%$ for all
five, four within $2\%$. Only the constant moves with the convention; the correlation does
not. That is what one wants: a change of definition rescales the fitted length without
destroying the relation.

\section{What this is not}\label{sec:limits}

\paragraph{The radius is supplied, not computed.} Equation~\eqref{eq:oneconst} relates two
measured quantities --- an ionization energy and an ionic radius --- through a function the
model provides. Given $r_{\rm ion}$, the model returns the ionization energy to a few per
cent. It does not produce $r_{\rm ion}$. The claim is therefore weaker than a prediction from
$Z$ alone, and stronger than a fit: a fit to five observables with five parameters asserts
nothing, whereas one constant reproducing five independent measurements can fail and does
not.

\paragraph{The model's own interface is not the screening radius.} It is natural to hope the
radius could be read off the atom instead of from a crystal. It cannot. A space-filling
radial partition locates a core--valence interface, and that interface is a property of
homogenising each shell onto a sphere; the screening radius is where the enclosed charge
leaves the valence seeing a net $+1$. The two do not track. Taken from the interface, the
radius gives lithium $0.1968$ against $0.1981$ but sodium only half its value, the interface
having grown to $5.38$~a.u.\ where the screening radius must be near $2.1$. Across one to
five core shells the interfaces come out at $1.90$, $4.01$, $5.38$, $9.18$ and
$9.42$~a.u.\ while physical core radii saturate near $2$--$3$: a partition giving every shell
a nearly fixed radial width, rather than contracting it with the charge that shell sees,
misplaces every interface past the first. Lithium works because its core \emph{is} the
innermost shell, where homogenisation barely moves anything.

\paragraph{Gold does not obey it.} Gold's measured first ionization energy, $0.339$~Ha, is
anomalously large, and the construction requires $R\approx0.18$~a.u.\ to reproduce it --- of
order $0.07\,r_{\rm ion}(\mathrm{Au}^+)$, not $0.44$. The relation
of~\eqref{eq:oneconst} is a statement about the alkali column and is not claimed beyond it.
What the construction does reproduce qualitatively is the \emph{direction}: a valence electron
outside an $18{+}1$ ending sits at a smaller radius than one outside $8{+}1$, hence binds
more tightly, which is the observed contrast between gold at $0.339$ and caesium at
$0.143$~\cite{NIST}.

\paragraph{An all-electron route exists and needs no radius at all.} Computed in three
dimensions, with the valence given angular freedom over homogenised inner shells and the same
$\Delta$SCF difference throughout, the first ionization energies of the first ten elements
come out with the correct sign everywhere and a mean absolute error of $18\%$, lithium and
sodium within $1\%$ and neon worst at $1.51\times$~\cite{Johnson-blog3D}. That route carries
no parameter and is the one to prefer where it is affordable; the construction of
Section~\ref{sec:constr} is a valence-only reduction, cheaper and with a supplied radius.
Computed radially instead --- the valence homogenised onto a sphere --- the same difference
gives carbon $0.26\times$ the measured value and the valence unbinds entirely by
$Z\approx14$. The failure is in the loss of the \emph{valence's} angular freedom, not in the
difference: carbon is $0.26\times$ radially and $0.97\times$ with the valence free, from the
same configuration.

\section{Conclusion}

Ionization energy cannot be obtained by differencing total energies at per-cent accuracy,
because from boron onward it is smaller than the error bar on the totals. A construction that
removes the core analytically rather than numerically avoids the difficulty, is exact at the
two points where it has no freedom, and --- once the spurious outer cutoff is removed, leaving
the inner radius as the only length --- reproduces the whole alkali column from a single
constant times a measured ionic radius, to four per cent. The radius comes from outside the
model, and the relation does not extend to gold. Within those limits it is a one-parameter
account of five independent measurements, and the kind of claim that can be refuted.

\begin{thebibliography}{99}

\bibitem{NIST} A.~Kramida, Yu.~Ralchenko, J.~Reader and NIST ASD Team, \emph{NIST Atomic
Spectra Database}, National Institute of Standards and Technology, Gaithersburg MD.

\bibitem{Shannon1976} R.~D.~Shannon, \emph{Revised effective ionic radii and systematic
studies of interatomic distances in halides and chalcogenides}, Acta Crystallogr.\ A
\textbf{32} (1976) 751--767.

\bibitem{Koopmans} T.~Koopmans, \emph{\"Uber die Zuordnung von Wellenfunktionen und
Eigenwerten zu den einzelnen Elektronen eines Atoms}, Physica \textbf{1} (1934) 104--113.

\bibitem{SzaboOstlund} A.~Szabo and N.~S.~Ostlund, \emph{Modern Quantum Chemistry},
Dover, 1996.

\bibitem{Hedin} L.~Hedin, \emph{New method for calculating the one-particle Green's function
with application to the electron-gas problem}, Phys.\ Rev.\ \textbf{139} (1965) A796.

\bibitem{Johnson-atoms} C.~Johnson, \emph{A 3D Multiphase Continuum Model for Atoms and the
Periodic Table}, submitted.

\bibitem{Johnson-blogAE} C.~Johnson, \emph{RealQM 1st ionization energy of alkaline earth
metals}, CJ on Mathematics and Science, 12 March 2025.
\url{https://claesjohnson.blogspot.com/2025/03/realqm-1st-ionization-energy-of.html}

\bibitem{Johnson-blog3D} C.~Johnson, \emph{Ionization: helium to neon}, CJ on Mathematics and
Science, 9 November 2023.
\url{https://claesjohnson.blogspot.com/2023/11/ionization-helium-to-fluorine.html}

\bibitem{Johnson-code1} C.~Johnson, \emph{Ionization 1 Valence New25}, p5.js sketch, 2025.
\url{https://editor.p5js.org/claesjohnson/full/gHG819EHB} --- the one-valence construction
of Section~\ref{sec:constr}, confining the density to the annulus $R<r<5$~a.u.

\bibitem{Johnson-code2} C.~Johnson, \emph{Ionization 2 Valence 1 Shells 0} and
\emph{\ldots Shells+1}, p5.js sketches, 2025.
\url{https://editor.p5js.org/claesjohnson/full/uAEDYKfbI},
\url{https://editor.p5js.org/claesjohnson/full/wGBhDVPT4E} --- the two-valence case, two
hemispherical domains on a kernel of $+2$, differenced as in~\eqref{eq:ie}.

\end{thebibliography}
\end{document}
